\section{The LongVALE Benchmark}
To build LongVALE, we propose an efficient and scalable pipeline that includes high-quality multi-modal long video filtering (Sec.~\ref{sec:filter}), omni-modal event boundary detection (Sec.~\ref{sec:boundary}), and omni-modal event captioning with audio-visual correlation reasoning (Sec.~\ref{sec:3.3}). 
The annotation process is illustrated in Fig.~\ref{fig:dataset_pipeline}.
{\color{black}More details are in Appendix.}


\subsection{Data Collection and Filtering}
\label{sec:filter}
We source videos from ACAV-100M~\cite{lee2021acav100m}, which contains video clips with high audio-visual correspondence, covering a wide variety of topics. 
We download raw videos on YouTube without cutting to maintain the integrity of the video content, where the videos span 30 seconds to 10 minutes. 
Then, we design a filtering strategy to obtain high-quality videos containing rich visual and audio semantics, as well as temporal dynamic information. 

Firstly, we use metadata to filter out low-quality videos with resolutions below 360p and retain only those with English transcripts.
To collect videos with diverse sounds (\eg, \textit{people clapping/laughing}, \textit{dog barking}), we exclude those where speech dominates, defined as when subtitles cover over 95\% of the duration.
Further, we remove videos with static content (\ie, slide shows) using PySceneDetect~\cite{Pyscenedetect} to detect scenes.
If a scene's average frame difference is below a threshold, it is considered static, and videos with over 80\% static scenes are filtered to ensure diverse motion content.
Finally, we select videos with segments that have consistent audio-visual semantics. 
We split each video into 5-second visual and audio segments, then use C-MCR~\cite{wang2023connecting}, an audio-visual contrastive learning model, to compute similarity scores between segment embeddings. Videos with at least one segment having audio-visual similarity above 0.25 are retained.
This effectively filters out videos with irrelevant audio-visual signals, such as edited or overdubbed audio (\eg, background music and narration).

As a result, from a total of 100K raw videos collected from ACAV-100M, we obtained 8.4K videos, highlighting our high standard for video quality.

\subsection{Omni-Modal Event Boundary Detection}
\label{sec:boundary}
Existing video benchmarks~\cite{wang2023internvid,chen2024panda,krishna2017dense,gao2017tall} segment videos into events based solely on visual signals. 
However, audio is equally crucial for a complete video content understanding, and audio events often have boundaries that do not align with visual ones. 
For example, a scene shift from a stage performance to the audience or a news broadcast cutting from a live scene to the studio may disrupt the visual event, but the audio continues.
Thus, relying only on visual cues for event boundary detection can severely break the semantic coherence of the audio.
To solve the issue, we propose to detect omni-modal event boundaries for the first time, considering both visual and audio scene boundaries. 

\noindent \textbf{Visual event boundary.}
Using only visual cues, we apply a two-stage detection method~\cite{chen2024panda} which includes splitting basic visual scenes and then merging semantically similar ones.
Notably, we refine the previous method to handle both very short and long visual events (2 seconds to 10 minutes) in long videos.
Besides, post-processing is also applied to exclude static scenes and transition clips, ensuring each visual event contains rich and meaningful content.

\noindent \textbf{Generic semantics-aware audio event boundary.}
Although audio is crucial for temporal video understanding, no method exists for detecting generic audio event boundaries without pre-defined categories.
To fill this gap, we design a generic method that segments long audio sequences into semantically coherent clips leveraging distinct audio properties.
Specifically, we first extract audio features using Mel-Frequency Cepstral Coefficients (MFCC)~\cite{juvela2018speech,hansen2015speaker},
which captures the audio's key frequency characteristics aligned with human auditory perception.
Then, we compute the mean of MFCC deltas (first-order differences) to summarize temporal variations, identifying values above a threshold as significant audio transitions.
Since such splitting only considers changes in spectral properties and overlooks semantic transitions, we further adopt CLAP~\cite{wu2023large} to extract audio embeddings, stitching adjacent clips if they are semantically similar.
We also implement additional procedures to merge segments split by abrupt changes, \eg, pauses between spoken words or sudden shifts in music.  

\noindent \textbf{Omni-modal event boundary.}
After identifying the single-modality event boundaries, we found visual events tend to be longer than audio ones, with multiple audio events often occurring within a single visual event. 
To define omni-modal event boundaries, we primarily rely on visual boundaries while preserving the integrity of audio events. 
Specifically, for each visual event, we set its start time as the beginning of the omni-modal event and include all overlapping audio events without truncation to determine the event's end. 
By doing this, we can maximize the semantic integrity and coherence of events across both modalities.


\subsection{Omni-Modal Event Captioning}
\label{sec:3.3}
We develop a comprehensive relation-aware captioning strategy to generate high-quality omni-modality dense captions for long videos by integrating all modality information (\ie, visual, audio, and speech). 
As shown in Fig.~\ref{fig:dataset_pipeline}, we first generate detailed dense captions for each modality, and then integrate them to obtain omni-modality captions, emphasizing the reasoning of semantic and temporal relationships across events from various modalities.


\begin{figure*}[ht]
  \centering
  \setlength{\abovecaptionskip}{1.0mm}
   \includegraphics[width=1.0\linewidth]{figures/new/fig3.jpg}
   \vspace{-5mm}
   \caption{Statistics of LongVALE benchmark. (a) Video duration distribution of both training and test sets. (b) Distribution of the number of omni-modal events in videos for both training and test sets. (c) Distribution of omni-modal event duration. (d) Distribution of audio-visual correlation types. The examples of omni-modal events with different audio-visual correlations are also illustrated.}
   \label{fig:dataset_overall}
     \vspace{-5mm}
\end{figure*}

\noindent \textbf{Dual-focus visual captioning.}
Existing video automatic captioning strategies~\cite{rizve2024vidla,chen2023vast,chen2024sharegpt4video} only use image captioners like BLIP~\cite{li2022blip} and GPT-4V~\cite{Gpt-4v} to describe uniformly sampled frames, lacking the awareness of temporal dynamic knowledge in videos. 
In contrast, we focus on both spatial visual details and significant dynamic information (\ie, actions and camera movements), and also consider the complexities of long video events. 
Specifically, we employ LLaVA-NeXT-Video~\cite{zhang2024llavanextvideo} to caption each video event. 
However, the model's performance drops with longer clips, so we divide clips longer than 30 seconds and caption them separately to preserve as many dynamic details as possible.
Additionally, we sample the center frame of each cropped clip as the keyframe and use GPT-4o~\cite{Gpt-4o} to generate comprehensive image captions. 
Incorporating such precise spatial details (\ie, OCR, object appearance, and scene context) helps improve caption quality and effectively correct errors caused by hallucinations from the video model.

\noindent \textbf{Audio and speech captioning.}
To capture the full details for each audio event, we focus on both general audio and speech captioning, as they are equally essential for audio content understanding.
Specifically, we employ Qwen-Audio~\cite{chu2023qwen}, a general large audio-language model, to obtain detailed audio descriptions. 
At the same time, Whisper-Large~\cite{radford2023robust}, a strong automatic speech recognition (ASR) model, is applied to get accurate subtitles if the clip contains speech content.  
Additionally, we perform further data refinement to minimize the impact of hallucinations. 


\noindent \textbf{Relation-aware omni-modal event captioning.}
Our goal is to generate high-quality omni-modal event captions by reasoning about audio-visual correlations and temporal dynamics across modality-specific events.
Specifically, instead of simply concatenating modality-specific captions~\cite{chen2023vast,liu2024valor}, we instruct Gemini-1.5-Pro~\cite{Gemini} to establish meaningful connections between them, such as analyzing whether audio events are visible, identifying sound sources, and reasoning about causality, \etc.
Additionally, we provide the model with single-modal event boundaries, guiding it to perceive fine-grained temporal changes within an omni-modal event, such as camera movements or sequential speaking in conversations, and summarize these to create fine-grained time-aware descriptions.  
Moreover, we feed the generated captions of previous omni-modal events as context, enabling the model to produce a more coherent and accurate caption for the current omni-modal event, thereby ensuring a more seamless dense narrative for a long video. 


\subsection{Subset Split and Manual Efforts}
\vspace{-1mm}
Through the above carefully designed scalable pipeline, we automatically collect 8,411 high-quality long videos with highly valuable dense omni-modal event annotations. 
We then meticulously divide the data into training and test sets, ensuring consistent data distribution (\eg, video duration and event counts) between them.
The final training and test sets consist of 7,240 and 1,171 long videos, respectively.

\noindent\textbf{Test set manual check and correction.}
To build a high-quality test set, we conducted thorough manual checks and corrections.
We developed an interface where one group reviews entire videos to verify the accuracy of omni-modal event boundaries and captions. Another group then corrects flagged errors to ensure high precision of annotations.



\subsection{Statistic Analysis}
\label{sec:3.5}
\vspace{-1mm}
Overall, our LongVALE is the first-ever omni-modality long video understanding benchmark with dense event-level annotations. The statistics are shown in Fig.~\ref{fig:dataset_overall}.
It includes 8,411 long videos spanning over 549 hours, with an average video duration of 3.9 minutes. 
Notably, the dataset contains 105,730 omni-modal events (91,863/13,867 in train/test split), each annotated with accurate temporal boundaries and omni-modality relation-aware captions.  

\noindent \textbf{Omni-modal event distribution.}
As shown in Fig.~\ref{fig:dataset_overall}(b), a large number of videos contain multiple omni-modal events, with an average of 12.6 events per video. 
Besides, the events have various lengths spanning from 1 second to even 10 minutes, with an average length of 16.7 seconds. 
Most events are relatively short, with 60\% lasting under 10 seconds and 97\% under 30 seconds as shown in Fig.~\ref{fig:dataset_overall}(c). 
Moreover, all events are non-overlapping and cover 89\% of the total video duration, highlighting the dense nature and complexity of real-world multi-modal long videos.

\noindent \textbf{Omni-modal caption characteristics.}
We highlight that omni-modal captions exhibit unique characteristics of \textit{audio-visual correlations} and \textit{fine-grained temporal dynamics}.
We further analyze all omni-modal captions in our test set using Gemini-1.5-Pro~\cite{Gemini} to identify these characteristics:
\textbf{1)} We define seven types of audio-visual correlations with two unimodal types (\ie, visual/audio-only). 
The occurrences of each type are shown in Fig.~\ref{fig:dataset_overall}(d).
It indicates that \textit{complementary}, \textit{synchronicity} and \textit{enhancement} are common correlation types. 
As shown in Fig.~\ref{fig:dataset_overall}, the speech content \textit{complements} the visual scene, where the visible woman speaker reflects \textit{synchronized} visual and audio events (example 1), and the off-screen commentator's excited tone \textit{enhances} the intense competition atmosphere (example 3).
Besides, there also exist other correlations, such as in the second example, where the crowd's cheers indicate a \textit{scene-aware} context of the stadium and the athlete's success triggers roars, demonstrating the \textit{causality} and \textit{temporal association} between the visual and audio events. 
Moreover, audio can provide \textit{corrective} insights for misleading visual cues, as seen in the fourth example where background laughter reveals a comedy show despite the man's serious face.
\textbf{2)} Additionally, the quantitative analysis also indicates that 78\% of the omni-modal captions capture fine-grained temporal dynamics. As highlighted in the blue words in Fig.~\ref{fig:dataset_overall}, the captions reflect a high-level understanding of sequential sub-events across different modalities, \eg, camera zooms and plot progressions. 
It underscores the extensive and complex semantic information embedded in our omni-modal captions.

